\chapter{Number Theoretic Transform (NTT)}
\label{ch:ntt}

\section{Purpose and scope}

The NTT is a finite-field analogue of a discrete Fourier transform. It changes representation so that structured polynomial multiplication becomes much cheaper. It is an integral part of ML-KEM's specified arithmetic, not merely an optional optimization~\cite{fips203}.

This chapter distinguishes three settings that are often conflated:
\begin{enumerate}
  \item a cyclic NTT for $\mathbb{F}_q[x]/(x^n-1)$;
  \item a full scalar negacyclic transform for $\mathbb{F}_q[x]/(x^n+1)$, when a primitive $2n$-th root exists; and
  \item ML-KEM's partial NTT, whose output consists of 128 degree-one residues modulo quadratic polynomials.
\end{enumerate}

\section{Roots of unity and the cyclic transform}

If $\omega\in\mathbb{F}_q$ has multiplicative order $n$, then evaluation at
\[
1,\omega,\omega^2,\ldots,\omega^{n-1}
\]
identifies multiplication modulo $x^n-1$ with $n$ scalar products. This requires $n\mid(q-1)$. The standard radix-2 butterfly network implements this map in $O(n\log n)$ operations when $n$ is a power of two.

For the negacyclic quotient $x^n+1$, scalar evaluation points must be roots of $x^n+1$, so a primitive $2n$-th root is required. The difference is structural, not cosmetic: an algorithm written only for an $n$-th root and cyclic convolution is not automatically an ML-KEM algorithm.

\subsection{Finding a root in a toy field}

The nonzero elements of $\mathbb{F}_{17}$ form a group of order $16$. If $g$ is a generator, then $g^2$ has order $8$ and can serve as $\omega$ for an 8-point cyclic NTT. SageMath can check that condition directly:

\begin{lstlisting}[language=Python]
F = GF(17)
g = F.multiplicative_generator()
omega = g^2
print(omega.multiplicative_order())  # 8
\end{lstlisting}

The important test is the multiplicative order, not merely that an element happens to satisfy an isolated power equation. For a coefficient vector $(a_0,\ldots,a_{n-1})$, the cyclic transform records
\[
\widehat a_i=a(\omega^i)=\sum_{j=0}^{n-1}a_j\omega^{ij},\qquad 0\le i<n.
\]
The inverse uses $\omega^{-1}$ and multiplies by $n^{-1}$, so an NTT is a reversible change of representation, not a lossy sampling operation.

\section{A complete toy cyclic NTT}

The following SageMath program performs a complete forward transform, pointwise product, and inverse transform in the cyclic ring $\mathbb{F}_{17}[x]/(x^8-1)$. It is intentionally a toy model, not ML-KEM.

\begin{lstlisting}[language=Python]
q = 17
F = GF(q)
n = 8
g = F.multiplicative_generator()
omega = g^((q - 1)//n)
assert omega.multiplicative_order() == n

def ntt(a):
    return [sum(a[j] * omega^(i*j) for j in range(n)) for i in range(n)]

def intt(A):
    n_inv = F(n)^(-1)
    return [n_inv * sum(A[i] * omega^(-i*j) for i in range(n))
            for j in range(n)]

a = vector(F, [1, 2, 3, 0, 0, 0, 0, 0])
b = vector(F, [4, 5, 0, 0, 0, 0, 0, 0])
A, B = ntt(a), ntt(b)
c = vector(F, intt([A[i]*B[i] for i in range(n)]))
print(c)                 # (4, 13, 5, 15, 0, 0, 0, 0)  (22 = 5 in GF(17))
print(intt(ntt(a)) == list(a))  # True
\end{lstlisting}

The code evaluates at powers of the actual primitive root \texttt{omega}, uses an inverse transform, and checks reconstruction. It therefore demonstrates an NTT, unlike evaluation at the unrelated points $0,1,\ldots,n-1$.

\subsection{Following one product through the transform}

For the input polynomials $a(x)=1+2x+3x^2$ and $b(x)=4+5x$, direct multiplication gives $4+13x+22x^2+15x^3$. The program computes the same result in four conceptual steps:
\begin{enumerate}
  \item evaluate $a$ at $1,\omega,\ldots,\omega^7$;
  \item evaluate $b$ at the same points;
  \item multiply the eight pairs of values; and
  \item apply the inverse transform to obtain $(4,13,22,15,0,0,0,0)$.
\end{enumerate}
Because the toy ring is modulo $x^8-1$ and the product has degree below eight, no visible wrap-around occurs. Higher-degree terms would wrap cyclically, not negacyclically.

\section{Butterflies and complexity}

The direct code is intentionally transparent, but it evaluates every output from every input and therefore takes $O(n^2)$ field operations. A practical radix-2 NTT uses butterflies to reuse work, by the same divide-and-conquer recursion Cooley and Tukey published for the DFT~\cite{cooleytukey1965}. A butterfly combines two values $u$ and $v$ using a twiddle factor $\omega^k$:
\[
(u,v)\longmapsto(u+\omega^k v,\;u-\omega^k v).
\]
There are $\log_2n$ stages and $O(n)$ work per stage, giving $O(n\log n)$ operations.

\begin{figure}[H]
\centering
\begin{tikzpicture}[x=1.2cm,y=0.72cm,every node/.style={font=\scriptsize}]
  \foreach \y/\name in {0/$a_0$,1/$a_1$,2/$a_2$,3/$a_3$}
    \node[left] at (0,-\y) {\name};
  \foreach \y/\name in {0/$\hat a_0$,1/$\hat a_1$,2/$\hat a_2$,3/$\hat a_3$}
    \node[right] at (4,-\y) {\name};
  \foreach \u/\v in {0/1,2/3} {
    \draw[pqcteal,thick] (0,-\u)--(2,-\u); \draw[pqcteal,thick] (0,-\u)--(2,-\v);
    \draw[pqcteal,thick] (0,-\v)--(2,-\u); \draw[pqcteal,thick] (0,-\v)--(2,-\v);
  }
  \foreach \u/\v in {0/2,1/3} {
    \draw[pqcblue,thick] (2,-\u)--(4,-\u); \draw[pqcblue,thick] (2,-\u)--(4,-\v);
    \draw[pqcblue,thick] (2,-\v)--(4,-\u); \draw[pqcblue,thick] (2,-\v)--(4,-\v);
  }
  \node[pqcteal] at (1,0.65) {stage 1};
  \node[pqcblue] at (3,0.65) {stage 2};
\end{tikzpicture}
\caption{A conceptual four-point radix-2 butterfly network. The picture shows the reuse pattern, not ML-KEM's complete partial NTT.}
\label{fig:toy-ntt-butterfly}
\end{figure}

\section{Why ML-KEM is different}

ML-KEM uses $q=3329$ and $n=256$. Since
\[
q-1=3328=2^8\cdot13,
\]
there are primitive $256$-th roots in $\mathbb{Z}_{3329}$, but no primitive $512$-th roots. In particular, $x^{256}+1$ does not split into 256 linear factors. With $\zeta=17$, a primitive 256-th root, FIPS~203 gives the factorization
\[
x^{256}+1=\prod_{i=0}^{127}\left(x^2-\zeta^{2\operatorname{BitRev}_7(i)+1}\right).
\]
Consequently, the ML-KEM NTT is an isomorphism from
\[
R_q=\mathbb{Z}_{3329}[x]/(x^{256}+1)
\]
to a product of 128 quadratic residue rings, not to 256 scalar evaluation values~\cite[Sec.~4.3]{fips203}.

\subsection{The three-level mental model}

\begin{table}[H]
\centering
\small
\begin{tabular}{p{0.24\textwidth}p{0.26\textwidth}p{0.31\textwidth}}
\toprule
Setting & Required roots / factorisation & Multiplication after transform \\
\midrule
Cyclic NTT, $x^n-1$ & primitive $n$-th root & $n$ scalar products \\
Full negacyclic NTT, $x^n+1$ & primitive $2n$-th root & $n$ scalar products at roots of $x^n+1$ \\
ML-KEM, $x^{256}+1$ over $\mathbb{Z}_{3329}$ & 128 quadratic factors & 128 quadratic base products \\
\bottomrule
\end{tabular}
\caption{Related transforms with different algebraic targets. Only the third row describes ML-KEM.}
\label{tab:ntt-three-settings}
\end{table}

\section{Multiplication in the ML-KEM NTT domain}

Each NTT-domain coordinate is represented by two coefficients, $a_0+a_1X$, reduced modulo $X^2-\gamma$. For two such residues, multiplication is
\[
(a_0+a_1X)(b_0+b_1X)\bmod(X^2-\gamma)
= (a_0b_0+a_1b_1\gamma)+(a_0b_1+a_1b_0)X.
\]
FIPS~203's \texttt{MultiplyNTTs} applies this base-case multiplication to each of the 128 coordinate pairs. Thus ``pointwise multiplication'' is accurate only if each point is understood as a quadratic component; it is not 256 independent scalar multiplications.

\begin{algorithm}[H]
\caption{ML-KEM NTT-domain multiplication (conceptual form)}
\label{alg:mlkem-base-multiply}
\begin{algorithmic}[1]
\Require NTT representations $\hat f,\hat g\in\mathbb{Z}_q^{256}$
\Ensure $\hat h=\hat f\times_{T_q}\hat g$
\For{$i=0$ to $127$}
  \State $\gamma\gets\zeta^{2\operatorname{BitRev}_7(i)+1}$
  \State $\hat h_{2i}\gets \hat f_{2i}\hat g_{2i}+\gamma\hat f_{2i+1}\hat g_{2i+1}\pmod q$
  \State $\hat h_{2i+1}\gets \hat f_{2i}\hat g_{2i+1}+\hat f_{2i+1}\hat g_{2i}\pmod q$
\EndFor
\State \Return $\hat h$
\end{algorithmic}
\end{algorithm}

The forward and inverse transforms themselves use butterfly networks and precomputed twiddle factors, as specified by FIPS~203. Their implementation details are deliberately not reconstructed here from a generic radix-2 routine: a conforming implementation must follow the standard's algorithms, bounds, and test vectors.

\section{Practical implementation checklist}

Before treating an NTT implementation as an ML-KEM implementation, verify all of the following.
\begin{itemize}
  \item The coefficient ring is $\mathbb{Z}_{3329}[x]/(x^{256}+1)$, not a cyclic substitute.
  \item The forward and inverse transforms follow the specified ordering and twiddle factors.
  \item NTT-domain multiplication operates on 128 pairs with the appropriate quadratic modulus.
  \item Modular reduction, serialization, sampling, and bounds follow FIPS~203 in addition to the transform itself.
  \item The implementation is checked against official known-answer tests and is engineered for its required side-channel model.
\end{itemize}

This checklist explains why a short classroom transform is valuable for intuition but must not be described as a complete implementation of ML-KEM.

\section{Summary}

An NTT is not simply ``evaluate a polynomial at several integers.'' It is a structured transform at roots of unity, together with its inverse. A cyclic transform uses $x^n-1$; a full scalar negacyclic transform needs a primitive $2n$-th root. ML-KEM uses neither simplified picture directly: its partial NTT maps $R_q$ to 128 quadratic components and multiplies them with base-case quadratic products.

\printbibliography[heading=chapterrefs]
